%% Chapter 7 — The Fourth Constraint: Extending the Universal Causal Equation

\chapter{The Fourth Constraint: Extending the Universal Causal Equation}
\label{chap:fourth-constraint}

\epigraph{%
  ``A chain is only as strong as its weakest link.''
}{--- proverb; in manufacturing intelligence, the weakest link was always the unverified reasoning engine}

\section{Recapitulation of the Prior Series}
\label{sec:prior-series}

Papers~2 and~3 of the \emph{wasm4pm} series established the Universal Causal Equation
in its original three-constraint form. We recall the relevant definitions here for
self-containment.

\paragraph{Constraint $C_0$ — Algorithmic Completeness.}
Every algorithm invoked by the manufacturing pipeline must be registered in the kernel
registry and reachable via the WASM boundary. A pipeline that silently falls back to an
unregistered algorithm violates $C_0$ and may not claim a valid receipt.

\paragraph{Constraint $C_1$ — Receipt Continuity.}
Every manufacturing stage must emit a BLAKE3 receipt whose \texttt{input\_hash} is the
\texttt{output\_hash} of the preceding stage. A gap or reorder in the receipt chain
breaks causal provenance and invalidates the pipeline's proof of correctness.

\paragraph{Constraint $C_2$ — Process Conformance.}
The object-centric event log derived from the pipeline's OTel spans must conform with
fitness $\phi \geq \phi^*$ (typically $\phi^* = 0.85$) to the declared process model
under pm4py token-based replay. A pipeline that passes code-level tests but fails
process conformance is non-conforming by the van der Aalst Constitution
(\texttt{\textasciitilde/.claude/rules/process-mining-chicago-tdd.md}).

Under $C_0 \wedge C_1 \wedge C_2$, the Universal Causal Equation reads:
\[
  A_U = \mu_U(\mathcal{O}^*_B)
  \quad \text{s.t.} \quad C_0 \wedge C_1 \wedge C_2.
\]

This formulation was sufficient for a pipeline whose reasoning engines were treated as
black boxes. Chapter~\ref{chap:bvc} reveals the gap: $C_0$--$C_2$ constrain the
\emph{pipeline} but not the \emph{breed} executing inside it. A breed that is
non-deterministic, or whose oracle is fakeable, or whose receipt does not cover the
ocel\_log field, may satisfy $C_0$--$C_2$ while producing reasoning artefacts that are
epistemically untrustworthy. The fourth constraint closes this gap.

\section{The Missing Constraint}
\label{sec:missing-constraint}

\begin{theorem}[Fourth Constraint Necessity]
\label{thm:fourth-constraint-necessity}
No finite set of pipeline-level constraints $\{C_0, C_1, C_2\}$ is sufficient to
guarantee trustworthy reasoning output if the participating breeds are uncertified.
\end{theorem}

\begin{proof}
We apply a No-Free-Lunch argument. Let $\mathcal{B}$ be any breed not required to
satisfy $C_{\mathrm{test}} = 1$. Then there exists an adversarial implementation
$\hat{\mathcal{B}}$ that:
(a) is registered in the kernel registry (satisfying $C_0$),
(b) emits a structurally valid BLAKE3 receipt chain (satisfying $C_1$),
(c) produces an OCEL log whose events have correct timestamps and object references,
    with fitness $\phi = 1.0$ under a trivially conforming single-trace model
    (satisfying $C_2$),
yet whose actual reasoning is the constant function $f_\perp(\cdot) = \perp$
(always returning a null or default conclusion regardless of input).

Concretely: $\hat{\mathcal{B}}$ could implement MYCIN by returning the first
diagnosis in alphabetical order irrespective of certainty factors. The pipeline
constraints $C_0$--$C_2$ cannot distinguish $\hat{\mathcal{B}}$ from a correct
implementation because they inspect only structural receipt topology and event log
conformance, not the epistemic content of the conclusion. The No-Free-Lunch theorem
for supervised learning (Wolpert 1996) generalises this: no finite set of structural
checks over the pipeline wrapper can certify the correctness of an arbitrary inner
function without a breed-specific oracle that constrains the inner function directly.
Therefore, a fourth constraint operating at the breed level is both necessary and
non-redundant with respect to $C_0$--$C_2$.
\end{proof}

\begin{definition}[The Fourth Constraint $C_3$]
\label{def:fourth-constraint}
The \emph{fourth constraint} $C_3$ requires that every breed $\Breed_i$ participating
in the universal manufacturing pipeline $\mu_U$ holds its Breed Validation Certificate:
\[
  C_3 \;:\; \forall\, \Breed_i \in \mu_U,\quad V(\Breed_i) = 1.
\]
Operationally, $C_3$ is enforced by the \texttt{admitBVC()} gate in
\texttt{packages/contracts/src/admission.ts}: at engine startup, \texttt{admitBVC()}
is called for every breed in the active registry. If any breed returns a
\texttt{BVCAdmissionResult} with \texttt{admitted = false}, the engine refuses to
start and emits a structured rejection with the failing component name and measured
value. This fail-fast behaviour ensures that $C_3$ cannot be bypassed by deferred
evaluation or lazy initialisation.
\end{definition}

\section{Theorem 7.2 --- The Four-Constraint Universal Equation}
\label{sec:four-constraint-equation}

\begin{theorem}[Four-Constraint Universal Causal Equation]
\label{thm:four-constraint-equation}
Under constraints $C_0 \wedge C_1 \wedge C_2 \wedge C_3$, the Universal Causal
Equation is:
\[
  \boxed{A_U \;=\; \mu_U\!\bigl(\mathcal{O}^*_B\bigr)
  \quad\text{s.t.}\quad
  C_0 \;\wedge\; C_1 \;\wedge\; C_2 \;\wedge\; C_3}
\]
where:
\begin{itemize}
  \item $A_U$ is the universal manufacturing artefact produced by the pipeline.
  \item $\mu_U$ is the universal manufacturing function.
  \item $\mathcal{O}^*_B$ is the object-centric process state enriched by the breed
        corpus $\mathcal{C}$.
  \item $C_0$ — algorithmic completeness (kernel registry coverage).
  \item $C_1$ — receipt continuity (BLAKE3 chain unbroken).
  \item $C_2$ — process conformance ($\phi \geq \phi^*$ under pm4py).
  \item $C_3$ — breed certification ($V(\Breed_i) = 1$ for all $\Breed_i \in \mu_U$).
\end{itemize}
\end{theorem}

\begin{proof}
By Theorem~\ref{thm:fourth-constraint-necessity}, $C_0 \wedge C_1 \wedge C_2$ is
insufficient. By Theorem~\ref{thm:universal-bvc}, all thirteen breeds currently satisfy
$C_3$. The conjunction $C_0 \wedge C_1 \wedge C_2 \wedge C_3$ is therefore both
necessary (removing any one constraint admits a counterexample by the argument of
Theorem~\ref{thm:fourth-constraint-necessity}) and sufficient (all four dimensions of
trustworthiness are covered): algorithmic completeness ensures no unregistered code runs;
receipt continuity ensures causal provenance is unbroken; process conformance ensures
the declared process model is respected; breed certification ensures each reasoning
engine is oracle-verified, receipt-chained, process-conformant at the breed level, and
deterministic. Together they constitute a closed proof obligation for a manufacturing
pipeline claiming to produce a certified artefact $A_U$.
\end{proof}

\section{Operational Semantics --- BVC as a Sheaf Condition}
\label{sec:sheaf-condition}

The four-constraint equation has a natural categorical reading that clarifies how local
breed certificates compose into a global pipeline certificate.

\begin{proposition}[BVC as a Sheaf Condition]
\label{prop:sheaf-condition}
Let the pipeline $\mu_U$ be modelled as a directed graph $G = (V, E)$ where each
node $v \in V$ corresponds to a manufacturing stage and each edge $e \in E$ carries a
receipt. Define a presheaf $\mathcal{F}$ on $G$ by assigning to each node $v$ the set
$\mathcal{F}(v) = \{A : \mathrm{BVC}(\Breed(v)) \wedge \text{receipt}(v)\}$ of
certified artefacts at that stage. Then $\mathcal{F}$ is a \emph{sheaf}: local
certificates glue to a global certificate if and only if $C_1$ (receipt continuity)
holds on all edges.
\end{proposition}

\begin{proof}
The sheaf axioms require (i) locality and (ii) gluing. Locality: if two stages $u, v$
agree on their shared receipt boundary (i.e., \texttt{output\_hash}$(u)$ =
\texttt{input\_hash}$(v)$, enforced by $C_1$), then the restriction of their
certified artefacts to the shared boundary is consistent. Gluing: given a compatible
family of locally certified artefacts at each node (each satisfying $C_3$), there
exists a unique globally certified artefact for the entire pipeline, namely $A_U$.
The uniqueness follows from the determinism component of $C_3$
($C_{\mathrm{determ}} = 1$): the same input seed always produces the same output,
so the global artefact is determined by the input and the pipeline topology, with no
ambiguity from non-deterministic choices at any stage.

To see that the sheaf condition fails without $C_3$: if any breed is non-deterministic
($C_{\mathrm{determ}} < 1$), then two pipeline runs with the same input may produce
different local sections at the non-deterministic node, and no global section can be
constructed from the resulting incompatible local data. The gluing axiom is thereby
violated. This confirms that $C_3$ is the algebraic completion of the sheaf structure
already implicit in $C_1$.

To see that the sheaf condition fails without $C_1$: if any edge lacks a valid receipt
(a gap in the BLAKE3 chain), the restriction maps are undefined on that edge, and
locality fails. Thus $C_1$ and $C_3$ are jointly necessary for the sheaf structure,
and together with $C_0$ and $C_2$ they provide the four independent proof obligations
that constitute the four-constraint universal equation.
\end{proof}

The sheaf reading has a practical consequence: the \texttt{admitBVC()} gate at engine
startup and the receipt chain validator at pipeline completion are not independent
checks. They are the local and global sections of the same sheaf, and verifying both
is equivalent to verifying the sheaf condition on the entire pipeline graph. An
audit tool that checks only one and not the other is examining a presheaf, not a
sheaf, and cannot issue a global conformance certificate.

\section{Summary}
\label{sec:fourth-constraint-summary}

This chapter completed the Universal Causal Equation by introducing $C_3$, the breed
certification constraint that operates at the level of individual reasoning engines
rather than the pipeline wrapper. Theorem~\ref{thm:fourth-constraint-necessity}
showed by a No-Free-Lunch argument that $C_0$--$C_2$ cannot detect an adversarial or
trivially-constant breed implementation, while Theorem~\ref{thm:four-constraint-equation}
established the complete four-constraint form $A_U = \mu_U(\mathcal{O}^*_B)$ s.t.\
$C_0 \wedge C_1 \wedge C_2 \wedge C_3$. Proposition~\ref{prop:sheaf-condition}
reframed the operational semantics as a sheaf condition, showing that determinism and
receipt continuity are two faces of the same categorical gluing requirement. The next
chapter turns to the latency and throughput properties of the certified breed corpus
under real workloads.
